\clearpage
\section{The Appointed Texts}
The Latin follows the Vatican typical edition of 1962. The English is the historical Lasance--Walsh witness, revised edition of 1945, in its registered later printing; it is study English, not an approved English text for celebration. Its printed errors are retained. The full source and boundary notes are in the terminal scope appendix.

\subsection{Introit}
\label{proper-introit}
\properlocus{1962 Missal: no. 3938, printed p. 697, left column; Lasance: printed 1257.}
\begin{latinproper}
Dignus est Agnus, qui occísus est, accípere virtútem, et divinitátem, et sapiéntiam, et fortitúdinem, et honórem. Ipsi glória et impérium in sǽcula sæculórum. \textnormal{Ps. 71, 1} Deus, iudícium tuum Regi da: et iustítiam tuam Fílio Regis. \textnormal{V.} Glória Patri, et Fílio, et Spirítui Sancto. \textnormal{R.} Sicut erat in princípio, et nunc, et semper, et in sǽcula sæculórum. Amen.

Dignus est Agnus, qui occísus est, accípere virtútem, et divinitátem, et sapiéntiam, et fortitúdinem, et honórem. Ipsi glória et impérium in sǽcula sæculórum.
\end{latinproper}
\begin{englishwitness}{Lasance--Walsh}
The Lamb that was slain is worthy to receive power, and divinity, and wisdom, and strength, and honor. To Him be glory and empire for ever and ever. \textnormal{Ps. 71, 1.} Give to the King Thy judgment, O God: and to the King's Son Thy justice. \textnormal{V.} Glory be to the Father, and to the Son, and to the Holy Ghost. \textnormal{R.} As it was in the beginning, is now, and ever shall be, world without end. Amen.

The Lamb that was slain is worthy to receive power, and divinity, and wisdom, and strength, and honor. To Him be glory and empire for ever and ever.
\end{englishwitness}
\subsection{Collect}
\label{proper-collect}
\properlocus{1962 Missal: no. 3939, printed p. 697, left column; Lasance: printed 1257.}
\begin{latinproper}
Omnípotens sempitérne Deus, qui in dilécto Fílio tuo, universórum Rege, ómnia instauráre voluísti: concéde propítius; ut cunctæ famíliæ géntium, peccáti vúlnere disgregátæ, eius suavíssimo subdántur império: Qui tecum vivit et regnat in unitáte Spíritus Sancti, Deus, per ómnia sǽcula sæculórum. \textnormal{R.} Amen.
\end{latinproper}
\begin{englishwitness}{Lasance--Walsh}
Almighty and eternal God, Who hast wished to restore all things through Thy beloved Son, the King of the universe, graciously grant that all the familes of the Gentiles separated by the wound of sin, may be subjected to His most loving dominion, Who with Thee liveth and reigneth in the unity of the Holy Ghost, God, world without end. \textnormal{R.} Amen.
\end{englishwitness}
\subsection{Epistle}
\label{proper-epistle}
\properlocus{1962 Missal: no. 3940, printed pp. 697–698; Lasance: printed 1257–1258.}
\begin{latinproper}
Léctio Epístolæ beáti Pauli Apóstoli ad Colossénses.

Fratres: Grátias ágimus Deo Patri, qui dignos nos fecit in partem sortis sanctórum in lúmine, qui erípuit nos de potestáte tenebrárum, et tránstulit in regnum Fílii dilectiónis suæ, in quo habémus redemptiónem per sánguinem eius remissiónem peccatórum: qui est imágo Dei invisíbilis, primogénitus omnis creatúræ: quóniam in ipso cóndita sunt univérsa in cælis et in terra, visibília et invisibília, sive throni, sive dominatiónes, sive principátus, sive potestátes: ómnia per ipsum et in ipso creáta sunt: et ipse est ante omnes, et ómnia in ipso constant. Et ipse est caput córporis Ecclésiæ, qui est princípium, primogénitus ex mórtuis: ut sit in ómnibus ipse primátum tenens: quia in ipso complácuit omnem plenitúdinem inhabitáre: et per eum reconciliáre ómnia in ipsum, pacíficans per sánguinem crucis eius, sive quæ in terris, sive quæ in cælis sunt, in Christo Iesu Dómino nostro.
\end{latinproper}
\begin{englishwitness}{Lasance--Walsh}
Lesson from the Epistle of blessed Paul the Apostle to the the Colossians.

Brethren, we give thanks to God the Father, Who hath made us worthy to be partakers of the lot of the saints in light: Who hath delivered us from the power of darkness, and hath translated us into the kingdom of the Son of His love. In Whom we have redemption through His blood, the remission of sins; Who is the image of the invisible God, the firstborn of every creature; for in Him were all things created in heaven and on earth, visible and invisible, whether thrones, or dominations, or principalities, or powers: all things were created by Him and in Him. And He is before all, and by Him all things consist. And He is the head of the body, the Church, Who is the beginning, the firstborn from the dead; that in all things He may hold the primacy: because in Him, it hath well pleased the Father, that all fulness should dwell; and through Him to reconcile all things unto Himself, making peace through the blood of His cross, both as to the things that are on earth, and the things that are in heaven, in Christ Jesus Our Lord.
\end{englishwitness}
\subsection{Gradual}
\label{proper-gradual}
\properlocus{1962 Missal: no. 3941, printed p. 698, left column; Lasance: printed 1258.}
\begin{latinproper}
Dominábitur a mari usque ad mare: et a flúmine usque ad términos orbis terrárum. \textnormal{V.} Et adorábunt eum omnes reges terræ: omnes gentes sérvient ei.
\end{latinproper}
\begin{englishwitness}{Lasance--Walsh}
He shall rule from sea to sea, and from the river unto the ends of the earth. \textnormal{V.} And all kings of the earth shall adore Him: all nations shall serve Him.
\end{englishwitness}
\subsection{Alleluia}
\label{proper-alleluia}
\properlocus{1962 Missal: no. 3942, printed p. 698, left column; Lasance: printed 1258–1259.}
\begin{latinproper}
Allelúia, allelúia. \textnormal{V.} \textnormal{Dan. 7, 14} Potéstas eius, potéstas ætérna, quæ non auferétur: et regnum eius, quod non corrumpétur. Allelúia.
\end{latinproper}
\begin{englishwitness}{Lasance--Walsh}
Alleluia, alleluia. \textnormal{V.} \textnormal{Dan. 7, 14.} His power is an everlasting power that shall not be taken away: and His kingdom that shall not be destroyed. Alleluia.
\end{englishwitness}
\subsection{Gospel}
\label{proper-gospel}
\properlocus{1962 Missal: no. 3945, printed p. 698; Lasance: printed 1259–1260.}
\begin{latinproper}
† Sequéntia sancti Evangélii secúndum Ioánnem.

In illo témpore: Dixit Pilátus ad Iesum: Tu es Rex Iudæórum? Respóndit Iesus: A temetípso hoc dicis, an álii dixérunt tibi de me? Respóndit Pilátus: Numquid ego Iudǽus sum? Gens tua et pontífices tradidérunt te mihi: quid fecísti? Respóndit Iesus: Regnum meum non est de hoc mundo. Si ex hoc mundo esset regnum meum, minístri mei útique decertárent, ut non tráderer Iudǽis: nunc autem regnum meum non est hinc. Dixit ítaque ei Pilátus: Ergo Rex es tu? Respóndit Iesus: Tu dicis, quia Rex sum ego. Ego in hoc natus sum, et ad hoc veni in mundum, ut testimónium perhíbeam veritáti: omnis qui est ex veritáte, audit vocem meam.
\end{latinproper}
\begin{englishwitness}{Lasance--Walsh}
† Continuation of the holy Gospel according to St. John.

At that time: Pilate said to Jesus: Art thou the King of the Jews? Jesus answered: Sayest thou this thing of thyself, or have others told it thee of Me? Pilate answered: Am I a Jew? Thy own nation, and the chief priests, have delivered thee up to me: what hast thou done? Jesus answered: My kingdom is not of this world. If My kingdom were of this world, My servants would certainly strive that I should not be delivered to the Jews: but now My kingdom is not from hence. Pilate therefore said to him: Art thou a king then? Jesus answered: Thou sayest that I am a King. For this was I born, and for this came I into the world; that I should give testimony to the truth, Everyone that is of the truth, heareth My voice.
\end{englishwitness}
\subsection{Offertory}
\label{proper-offertory}
\properlocus{1962 Missal: no. 3946, printed p. 698, right column; Lasance: printed 1260.}
\begin{latinproper}
Póstula a me, et dabo tibi gentes hereditátem tuam, et possessiónem tuam términos terræ.
\end{latinproper}
\begin{englishwitness}{Lasance--Walsh}
Ask of Me, and I will give The the Gentiles for Thy inheritance, and the utmost parts of the earth for Thy possession.
\end{englishwitness}
\subsection{Secret}
\label{proper-secret}
\properlocus{1962 Missal: no. 3947, printed p. 698, right column; Lasance: printed 1260.}
\begin{latinproper}
Hóstiam tibi, Dómine, humánæ reconciliatiónis offérimus: præsta, quǽsumus; ut, quem sacrifíciis præséntibus immolámus, ipse cunctis géntibus unitátis et pacis dona concédat, Iesus Christus, Fílius tuus, Dóminus noster: Qui tecum vivit et regnat in unitáte Spíritus Sancti, Deus, per ómnia sǽcula sæculórum. \textnormal{R.} Amen.
\end{latinproper}
\begin{englishwitness}{Lasance--Walsh}
O Lord, we offer Thee this host for the reconciliation of humanity; grant, we beseech Thee, that Jesus Christ Thy Son our Lord, Whom we immolate in this sacrifice, will bestow on all Gentiles the gifts of unity and peace, Who with Thee liveth and reigneth in the unity of the Holy Ghost, God, world without end. \textnormal{R.} Amen.
\end{englishwitness}
\subsection{Preface}
\label{proper-preface}
\properlocus{1962 Missal: no. 1080, printed p. 291; Lasance: printed 813–814, dialogue from 772.}
\begin{latinproper}
\textnormal{V.} Per ómnia sǽcula sæculórum. \textnormal{R.} Amen. \textnormal{V.} Dóminus vobíscum. \textnormal{R.} Et cum spíritu tuo. \textnormal{V.} Sursum corda. \textnormal{R.} Habémus ad Dóminum. \textnormal{V.} Grátias agámus Dómino Deo nostro. \textnormal{R.} Dignum et iustum est.

Vere dignum et iustum est, æquum et salutáre, nos tibi semper et ubíque grátias ágere: Dómine, sancte Pater, omnípotens ætérne Deus: Qui unigénitum Fílium tuum, Dóminum nostrum Iesum Christum, Sacerdótem ætérnum et universórum Regem, óleo exsultatiónis unxísti: ut, seípsum in ara Crucis hóstiam immaculátam et pacíficam ófferens, redemptiónis humánæ sacraménta perágeret: et, suo subiéctis império ómnibus creatúris, ætérnum et universále regnum imménsæ tuæ tráderet maiestáti: regnum veritátis et vitæ; regnum sanctitátis et grátiæ; regnum iustítiæ, amóris et pacis. Et ídeo cum Angelis et Archángelis, cum Thronis et Dominatiónibus, cumque omni milítia cæléstis exércitus, hymnum glóriæ tuæ cánimus, sine fine dicéntes:

Sanctus, Sanctus, Sanctus Dóminus Deus Sábaoth. Pleni sunt cæli et terra glória tua. Hosánna in excélsis. Benedíctus qui venit in nómine Dómini. Hosánna in excélsis.
\end{latinproper}
\begin{englishwitness}{Lasance--Walsh}
\textnormal{V.} World without end. \textnormal{R.} Amen. \textnormal{V.} The Lord be with you. \textnormal{R.} And with thy spirit. \textnormal{V.} Lift up your hearts. \textnormal{R.} We have them lifted up unto the Lord. \textnormal{V.} Let us give thanks to the Lord our God. \textnormal{R.} It is meet and just.

It is truly meet and just, right and profitable, for us, at all times, and in all places, to give thanks to Thee, O Lord, the holy One, the Father almighty, the everlasting God: Who didst anoint with the oil of gladness Thine only-begotten Son, our Lord Jesus Christ, eternal Priest and King of the universe: that, offering Himself a stainless peace-offering on the altar of the cross, He might fulfill the pledges of man's redemption; and, having all creatures subject to His power, might deliver to Thy sublime majesty an eternal and universal kingdom, a kingdom of truth and life; a kingdom of holiness and grace, a kingdom of justice, love and peace. And therefore with the angels and archangels, with the thrones and dominations, and with all the array of the heavenly host, we sing a hymn to Thy glory and unceasingly repeat:

Holy, Holy, Holy, Lord God of hosts. Heaven and earth are full of Thy glory. Hosanna in the highest. Blessed is He Who cometh in the name of the Lord. Hosanna in the highest.
\end{englishwitness}
\subsection{Communion}
\label{proper-communion}
\properlocus{1962 Missal: no. 3948, printed p. 698, right column; Lasance: printed 1260.}
\begin{latinproper}
Sedébit Dóminus Rex in ætérnum: Dóminus benedícet pópulo suo in pace.
\end{latinproper}
\begin{englishwitness}{Lasance--Walsh}
The Lord shall sit King forever. The Lord will bless His people with peace.
\end{englishwitness}
\clearpage
\subsection{Postcommunion}
\label{proper-postcommunion}
\properlocus{1962 Missal: no. 3949, printed p. 698, right column; Lasance: printed 1260.}
\begin{latinproper}
Immortalitátis alimóniam consecúti, quǽsumus, Dómine: ut, qui sub Christi Regis vexíllis militáre gloriámur, cum ipso, in cælésti sede, iúgiter regnáre possímus: Qui tecum vivit et regnat in unitáte Spíritus Sancti, Deus, per ómnia sǽcula sæculórum. \textnormal{R.} Amen.
\end{latinproper}
\begin{englishwitness}{Lasance--Walsh}
Fed with this immortal nourishment, we beseech Thee, O Lord, that we who glory to fight under the standard of Christ the King, may forever reign with Him on the heavenly throne. Who with Thee liveth and reigneth in the unity of the Holy Ghost, God, world without end. \textnormal{R.} Amen.
\end{englishwitness}
